\section{Machine-Checked Courseware}
\label{sec:course}

The course layer (\texttt{course/}) contains 22 modules
(Figure~\ref{fig:lifecycle}), each with the same anatomy, and is kept
consistent with the implementation by three guards:
\textbf{G1}~(\emph{equivalence-graded exercises}): every build exercise is
scored by functional equivalence against the maintained module it teaches;
\textbf{G2}~(\emph{drift-guarded exposition}): every identifier in displayed
code must exist in the referenced source file and in the page that shows it,
enforced by tests; \textbf{G3}~(\emph{claim--command coupling}): every
quantitative statement names the committed command that regenerates it
(Figure~\ref{fig:guards}). The paragraphs below describe the anatomy these
guards protect.

\begin{figure}[t]
  \centering
  \resizebox{\columnwidth}{!}{% Figure: the three machine-checked-courseware guards (Section 4)
\begin{tikzpicture}[
  font=\scriptsize,
  box/.style={draw, rounded corners=2pt, align=center, inner sep=2.5pt, fill=blue!6},
  ci/.style={draw, rounded corners=2pt, align=center, inner sep=2.5pt, fill=red!8},
  arr/.style={{Stealth[length=1.8mm]}-{Stealth[length=1.8mm]}, thick}
]
\node[box, text width=17mm, minimum height=9mm] (docs)  {exposition\\ \tiny displayed code};
\node[box, text width=17mm, minimum height=9mm, right=5.5mm of docs] (claims) {claims\\ \tiny measured numbers};
\node[box, text width=17mm, minimum height=9mm, right=5.5mm of claims] (labs) {exercises\\ \tiny learner code};
\node[box, minimum width=62mm, minimum height=7mm, below=9.5mm of claims] (impl)
  {\textbf{maintained implementation} \, {\tiny(\texttt{baby\_whale\_v4}, tested)}};
\draw[arr] (docs)   -- (docs |- impl.north)   node[font=\tiny, midway, left=0.5mm]  {G2};
\draw[arr] (claims) -- (claims |- impl.north) node[font=\tiny, midway, right=0.5mm] {G3};
\draw[arr] (labs)   -- (labs |- impl.north)   node[font=\tiny, midway, right=0.5mm] {G1};
\node[ci, right=3.5mm of impl, minimum height=7mm] (ci) {\tiny CI fails on\\ \tiny mismatch};
\draw[-{Stealth[length=1.8mm]}, thick, dashed] (impl) -- (ci);
\node[font=\tiny, below=0.5mm of impl.south west, anchor=north west]
  {G1 behavior $\equiv$ \quad G2 symbols exist \quad G3 regenerated by script};
\end{tikzpicture}
}
  \caption{\textbf{Machine-checked courseware.} Each artifact type is bound
  to the maintained implementation by a guard (G1 behavioral equivalence, G2
  symbol existence, G3 regeneration scripts); CI fails when any binding
  breaks.}
  \label{fig:guards}
\end{figure}

\paragraph{Five beats.} Every module reads as: \emph{the wall} (what breaks
without this feature), \emph{the idea} (theory and the paper), \emph{from
theory to code} (a term-by-term table mapping each piece of the math---or,
for systems modules, each rule of the mechanism---to the exact operation in
the source, with a ``why'' per row), \emph{the payoff, measured} (a command
that prints the number), and \emph{break it} (a knob to turn plus a
quantitative reflection question). Where an idea is genuinely mathematical the
table's left column is an equation typeset in \LaTeX{}; where it is a data
structure or protocol (paged KV, batching) the left column is the rule---the
bridge is the same.

\paragraph{Three lenses.} Modules are read through modeling/theory,
training/alignment, and systems lenses; a systems helper prints concrete
costs (parameters, bf16 vs.\ FP4 memory, active experts, attention MACs) for
any configuration.

\paragraph{Three tracks, graded by the test suite.}
\emph{Read} is the annotated narrative, which pastes the \emph{actual} key
lines of the implementation. \emph{Build} is a set of 23 fill-in-the-blank
labs---at least one per module---each a \texttt{NotImplementedError} stub
whose docstring carries the derivation; running the lab grades the learner's
code against the real module (a learner's transformer layer must match the
library's block tensor-for-tensor). The test suite doubles as the autograder,
so correctness cannot be faked. \emph{Extend} prompts are open problems that
become contributions.

\paragraph{Progressive motivation.} A preset ladder
(\texttt{gpt-minimal} $\to$ \texttt{+mla} $\to$ \texttt{+compressed} $\to$
\texttt{+moe-balanced} $\to$ \texttt{+mtp} $\to$ \texttt{full}) turns each
feature on one at a time, so learners hit the wall a feature removes before
enabling it.

\paragraph{Milestones.} Passing labs proves understanding of parts; five
milestones prove the assembled system works on real tasks (\emph{it learns};
\emph{it remembers}, via a long-context needle probe with per-sample varying
answers; \emph{it follows}; \emph{it reasons}, via pass@1; \emph{it serves}).
Two milestones train real (tiny) models inside the test suite and are gated
on every CI run.

\paragraph{Enforcement.} G1 is realized as one grader per exercise, each
itself tested to accept the reference solution and to reject a representative
wrong one. G2 is a snippet audit over all twenty content modules, paired with
a coverage check that no module lacks a guarded mapping or a graded exercise.
G3 is realized by committing the measurement scripts and their outputs
alongside the paper (\S\ref{sec:eval}). The net effect: renaming a single
function in the library fails the build until the courseware is
updated---documentation becomes falsifiable in the same sense the code is.
